\section{单调拉格朗日子流形的例子}
\label{sec:examples-monotone-lagrangians}

如二维例子所示，一般拉格朗日子流形对子 $(L_0,L_1)$ 的\personFloer{}同调会遇到障碍。
对这一障碍及其形变理论的完整研究需要超出本书层次的数学；(FOOO09) 已对此作了
彻底处理，有意了解细节的读者可参阅该书。

\personFloer{}在 (Fl89b) 中引入单调辛流形的概念，以克服其研究哈密顿\personFloer{}同调时遇到的
横截性困难。

\begin{remark}
\label{rem:general-perturbation-transversality}
在一般构造\personGromov{}--\personWitten{}不变量所涉及的模问题中，利用殆复结构以外的扰动研究
横截性的一般途径，最早由\personKontsevich{}在其著名论文 (Kon95) 中提出。

\stfhIndex{I117}{Cangxijiegou}{仓西结构（Kuranishi structure）}
\stfhIndex{I155}{Duozhongzhedietujing}{多重折叠途径（polyfold approach）}
如今，这种一般扰动横截性理论本质上有两种不同途径：一种是\emph{有限维壳上}途径，另一种
是\emph{无穷维离壳}途径。前一种是\personFukaya{}与\personOno{} (FOn99) 所采用的\termKuranishiStructure{}与多截面
方法，它是求解非线性方程的经典\personLyapunov{}约化方法之加强形式（相关构造亦见
(LT98)）；代数几何学家在 (LiT98)、(Be97)、(BF97) 中也采用这一途径。后一种途径
由\personBerndSiebert{} (Sie99) 采用，较近则以\personHofer{}、\personKrzysztofWysocki{}与\personZehnder{}的\termPolyfold{}途径
(HWZ07)、(H08) 为代表，尝试在离壳设置中对原方程作整体扰动（相关构造亦见
(Ru99)）。

即使 (FOn99) 问世二十年之后，这个横截性问题仍然是一个即便在辛几何专家之间也
需要进一步澄清的问题。
\end{remark}

受\personFloer{}对单调辛流形研究的启发，作者在 (Oh93a) 中挑出了单调拉格朗日子流形这一
类，以便把\personFloer{}对对子 $(L_0,L_1)$ 的同调理论构造推广到正合情形之外。这类拉格朗日
子流形避开了开\personRiemann{}曲面（亏格为零）的紧化模空间所涉及的大部分复杂横截性问题，
并使人能够无缝推广\personFloer{}的定义，依定义 16.1.1 中\personFloer{}原来的定义来定义映射
$\partial:CF(L_0,L_1)\to CF(L_0,L_1)$。

然而，(Oh93a) 已经注意到、(Oh95a) 后来又进一步强调，即使在这一良好情形中，也可能
出现\termAnomaly{}，即
\begin{equation}
 \partial^2\neq0,
\label{eq:monotone-boundary-anomaly}
\end{equation}
除非对子 $(L_0,L_1)$ 满足某个与马斯洛夫指标为 $2$ 的全纯圆盘之几何密切相关的
相容性条件。

\begin{remark}
\label{rem:role-of-monotonicity}
从内蕴辛拓扑的观点看，单调性既不自然也不灵活，因此它在一般辛拓扑中的作用对作者
而言尚未完全清楚。不过，它足以丰富地展示全纯圆盘如何在\personFloer{}同调中产生\emph{\termAnomaly{}}与
\emph{障碍}，并在\emph{非正合}背景下检验各种\personFloer{}理论构造。以某种方式处理上述障碍后，它也
自然推广了曲面上拉格朗日子流形\personFloer{}同调的讨论。在尝试借助障碍--形变理论与\termKuranishiStructure{}技术，把各种代数和几何结构完整实现于一般拉格朗日子流形的背景中
之前，单调情形是一个很好的试验场。近来，与（几乎）环面流形及镜像对称相关的许多
有趣单调拉格朗日环面例子已被发现；见 (FOOO12a)、(Wu12)、(Via13, Via14)。
\end{remark}

下面给出单个拉格朗日子流形版本的 $G(L_0,L_1;\ell_0)$；它也类似于环境辛流形
$(M,\omega)$ 的群
\[
 \Gamma_\omega=\Gamma(M,\omega)
 =\frac{\pi_2(M)}{\ker\omega\cap\ker c_1}.
\]
以 $I_\omega:\pi_2(M,L)\to\mathbb R$ 表示辛面积的取值。另定义整值同态
$I_\mu:\pi_2(M,L)\to\mathbb Z$：
\[
 I_\mu(\beta)=\mu\bigl(w^*TM,(\partial w)^*TL\bigr),
\]
其中右端是丛对
\[
 (w^*TM,(\partial w)^*TL)
\]
的马斯洛夫指标，而 $w:(D^2,\partial D^2)\to(M,L)$ 是 $\beta$ 的一个（因而
任意一个）代表。

\begin{definition}
\label{def:lagrangian-novikov-group}
定义
\[
 G(L;\omega)=\frac{\pi_2(M,L)}{\ker I_\omega\cap\ker I_\mu},
\]
并把 $\Lambda_{(L;\omega)}$ 定义为相应的\personNovikov{}环。
\end{definition}

简要回顾\personNovikov{}环 $\Lambda_{(\omega,L)}(R)$ 及其子环
$\Lambda_{0,(\omega,L)}(R)$ 的基本性质，其中 $R$ 是交换环，例如可取
$\mathbb Z_2$、$\mathbb Z$ 或 $\mathbb Q$。令
\[
 q^\beta=T^{\omega(\beta)}e^{\mu_L(\beta)},
\]
并令
\[
 \deg(q^\beta)=\mu_L(\beta),\qquad E(q^\beta)=\omega(\beta),
\]
于是 $\Lambda_{(\omega,L)}$ 与 $\Lambda_{0,(\omega,L)}$ 一般成为分次环。典范赋值
$\nu:\Lambda_{(\omega,L)}\to\mathbb R$ 定义为
\[
 \nu\left(\sum_\beta a_\beta T^{\omega(\beta)}e^{\mu_L(\beta)}\right)
 =\min\{\omega(\beta)\mid a_\beta\neq0\}.
\]
它诱导子环 $\Lambda_{0,(\omega,L)}\subset\Lambda_{(\omega,L)}$ 上的赋值，并进而
诱导其上的自然滤过。这使 $\Lambda_{(\omega,L)}$ 成为滤过分次环。对一般拉格朗日
子流形，该环甚至不必是\personNoether{}环。

下面定义单调拉格朗日子流形。

\stfhIndex{I131}{JixiaoMasiluofushu}{极小马斯洛夫数（minimal Maslov number）}
\stfhSubIndex{I133}{Dandiaolagelangriziliuxing}{单调拉格朗日子流形（monotone Lagrangian submanifold）}{Dingyi}{定义（definition）}
\begin{definition}
\label{def:monotone-lagrangian}
若紧拉格朗日子流形 $L\subset(M,\omega)$ 满足
\[
 I_\omega=\lambda I_\mu\qquad\text{对某个 }\lambda\geq0,
\]
则称 $L$ 是单调的。定义\emph{\termMinimalMaslovNumber{}} $\Sigma_L$ 为
$\operatorname{Im}I_\mu$ 的正生成元。
\end{definition}

\begin{exercise}
\label{ex:monotone-lagrangian-ambient}
证明单调拉格朗日子流形只存在于单调辛流形上。
\end{exercise}

\stfhIndex{I196}{Ruozhenghe}{弱正合（weakly exact）}
注意，所谓弱正合情形，即 $I_\omega=0$ 的情形，平凡地包含在上述情形中。若忽略
分次问题，则它可以完全按照正合情形 $\pi_2(M,L)=\{0\}$ 来研究。因此，本小节讨论
中假设 $L\subset M$ 单调、$\lambda>0$ 且 $I_\mu\neq0$。

下面的命题展示了这类拉格朗日子流形相对于一般拉格朗日子流形所具有的某些特殊辛
拓扑性质。

\begin{proposition}
\label{prop:monotone-lagrangian-group-rank-one}
设 $L\subset(M,\omega)$ 是单调拉格朗日子流形。则 $G(L;M,\omega)$ 是秩为 $1$
的自由\personAbel{}群，而 $\Lambda_{0,(\omega,L)}$ 成为一个离散赋值环，其剩余域是
$\Lambda_{(\omega,L)}$。
\footnote{译者注：原书此前把该群定义为 $G(L;\omega)$，本命题改写为
$G(L;M,\omega)$；且把 $\Lambda_{(\omega,L)}$ 称作
$\Lambda_{0,(\omega,L)}$ 的“剩余域”。译文照录，留待统一裁决。}
\end{proposition}

\stfhIndex{I056}{Kelifudehuanmian}{克利福德环面（Clifford torus）}
$\mathbb{RP}^n\subset\mathbb{CP}^n$ 以及例 5.19 中的\personClifford{}环面，都是单调
拉格朗日子流形的例子。

\begin{proposition}
\label{prop:clifford-torus-monotone}
考虑\personClifford{}环面 $T^n\subset\mathbb P^n$：
\[
 T^n=\{[z_0:z_1:\cdots:z_n]\in\mathbb P^n\mid |z_i|=1\}.
\]
该环面单调，其相伴群 $G(T^n;\mathbb P^n)$ 同构于 $\mathbb Z$。
\end{proposition}

\begin{proof}
容易验证
\[
 \pi_2(\mathbb P^n,T^n)\cong\pi_1(T^n)\oplus\pi_2(\mathbb P^n).
\]
若以 $\beta_i$ 表示圆盘
\[
 D_i=\{[z_0:z_1:\cdots:z_n]\in\mathbb P^n
 \mid z_j\equiv e^{\sqrt{-1}\theta_i^0}\text{ 对 }j\neq i\text{ 固定}\}
\]
之类，则在 $\pi_2(\mathbb P^n,T^n)$ 中有关系
\[
 \beta_0+\cdots+\beta_n=\alpha,
\]
其中 $\alpha$ 是 $\pi_2(\mathbb P^n)$ 的（正）生成元。又有
\begin{equation}
\begin{aligned}
 \mu(\beta_i)&=2=\frac{c_1(\alpha)}{n+1},\\
 \omega(\beta_i)&=\frac{\omega(\alpha)}{n+1}=2\pi.
 \end{aligned}
\label{eq:clifford-torus-monotonicity}
\end{equation}
\footnote{译者注：对 $\mathbb P^n$ 的标准正生成元通常有
$c_1(\alpha)=n+1$，故原书所印 $2=c_1(\alpha)/(n+1)$ 与此不相容；译文照录。}
这就证明了单调性。

此外，\eqref{eq:clifford-torus-monotonicity} 还蕴含所有 $\beta_i$ 在
$G(T^n;\mathbb P^n)$ 中定义同一个元素，记为 $\beta$，且
$\alpha=(n+1)\beta$。因此
\[
 G(T^n;\mathbb P^n)=\mathbb Z\cdot[\beta]\cong\mathbb Z.
\]
\end{proof}

再考虑其他标准环面
\[
 T_{(c_0,\ldots,c_n)}
 =\{[z_0:\cdots:z_n]\in\mathbb P^n\mid |z_i|=c_i>0\}.
\]
它们并不单调。以 $\beta_i$ 和 $\alpha$ 表示与上述
$c_0=\cdots=c_n$ 情形类似定义的类。依赖于数集
\[
 \{\omega(\beta_0),\ldots,\omega(\beta_n),\omega(\alpha)=2\pi\}
\]
的有理相关性，$G(T_{(c_0,\ldots,c_n)};\mathbb P^n)$ 的秩在 $1$ 与 $n+1$
之间变化。该例可进一步推广到环面流形的拉格朗日环面纤维；(CO06)、
(FOOO10b, FOOO11a) 对此作了细致分析。

单调性条件主要是为寻找一个条件而设计的，使 (Oh93a) 所定义的朴素\personFloer{}同调能够
构造。不过，第 3.7.2 节所述拉格朗日悬挂构造保持这一条件。

\begin{proposition}
\label{prop:monotone-suspension}
设 $L\subset(M,\omega)$ 是任意紧单调拉格朗日嵌入，$(h,\psi)$ 是第 3.7.2 节
意义下它的正合拉格朗日圈。则相伴悬挂
\[
 \iota_{(h,\psi)}:S^1\times L\longrightarrow T^*S^1\times M
\]
关于 $T^*S^1\times M$ 上的直和辛形式 $\omega_0\oplus w$ 是具有相同\termMonotonicityConstant{}的
单调拉格朗日嵌入。
\footnote{译者注：原书此处的直和辛形式印作 $\omega_0\oplus w$，而环境辛形式一直
记作 $\omega$；译文照录。}
\end{proposition}

\begin{exercise}
\label{ex:prove-monotone-suspension}
证明该命题。
\end{exercise}

还从命题 3.7.7 回忆，$\iota_{(h,\psi)}$ 在 $T^*S^1\times M$ 中哈密顿同痕于乘积
嵌入 $o_{S^1}\times L$。

\stfhIndex{I154}{Boerjieluoweiqi}{波尔捷罗维奇（Polterovich）}
\personPolterovich{}在研究 $\Ham(S^2)$ 的\personHofer{}直径时利用了上述命题，并把拉格朗日\personFloer{}同调的\personKunneth{}公式的一个版本应用于 $S^2$ 的赤道，巧妙地证明了下面的定理。第
17.6 节将依原论文 (Po98b) 给出证明。

\begin{theorem}[\personPolterovich{}]
\label{thm:hofer-diameter-s2-infinite}
$\Ham(S^2)$ 的\personHofer{}直径为无穷。
\end{theorem}

下面说明第 3.7.3 节所述\personChekanov{}环面 (Che96b) 的一个变体。

\begin{example}[在 $S^2\times S^2$ 中的\personChekanov{}环面]
\label{ex:chekanov-torus-s2xs2}
\personChekanov{}的环面最初通过把第 3.7.3 节所述悬挂构造用于标准圆周
$S^1\subset\mathbb C$ 而在 $\mathbb C^2$ 中构造（见例 16.3.6(3)）。适当选择\personChekanov{}悬挂中出现的参数，可将该环面嵌入 $S^2\times S^2$，成为一个奇异单调
拉格朗日环面：它与乘积
$S^1_{\mathrm{eq}}\times S^1_{\mathrm{eq}}\subset S^2\times S^2$ 不哈密顿同痕
（证明见 (EP10)、(FOOO12a)）。其构造可描述如下（此描述见 (AF08)）。把
$S^2\times S^2\setminus\overline\Delta$ 等同于单位圆盘丛 $D^1(TS^2)$。
(AF08) 中研究的拉格朗日环面，是 $\partial D^r(TS^2)$ 上经过一个给定点（比如北极
$N$）的简单闭定向测地线之并。适当选取 $r>0$，可以证明该环面在
$S^2\times S^2\supset D^1(TS^2)$ 中单调。把它记作 $T(r)$。
\end{example}

\begin{exercise}
\label{ex:unique-monotone-chekanov-radius}
证明上述例子中恰有一个 $0<r<1$ 的选择使环面 $T(r)$ 单调，并求出 $r$。
\end{exercise}

\personChekanov{}与\personSchlenk{} (CS10) 在辛向量空间或复射影空间中构造了许多单调拉格朗日环面，
并研究了其中哪些例子彼此不同。
